\subsection{Case Study: ChatBDI}
\label{sec:chatbdi}
\slidesrc{Terminology introduction: 11--22}

The course shows a real example from the lecturer's recent research, at the
intersection of multi-agent systems (S\&DAI) and NLP. It is useful early on because it
shows, on a concrete problem, almost every topic of the course: syntax, semantics,
pragmatics, edit distance and LLMs.

\subsubsection{The idea: think BDI, talk LLM}
\slidesrc{Terminology introduction: 12--13}

A \textbf{BDI agent} (Belief--Desire--Intention) is a software agent that reasons
symbolically: it has \emph{beliefs} about the world (facts such as
\texttt{has(owner, beer)}), \emph{desires} (goals) and \emph{intentions} (plans it has
committed to). BDI agents in a multi-agent system (MAS) talk to each other with an
\textbf{Agent Communication Language (ACL)}: each message has a
\textbf{performative} (the type of communicative act, e.g.\ \texttt{tell} to inform,
\texttt{askOne}/\texttt{askAll} to ask a question, \texttt{achieve} to ask for a goal
to be pursued) and a \textbf{content}, a logical term.

ChatBDI keeps the BDI agent as the ``brain'' and uses an LLM only as its
\textbf{language sensor--actuator}: the LLM translates a human's natural-language sentence
into an ACL message the agent understands (and the agent's answers back into natural
language). So humans can chat with agents that were never designed to understand
natural language.

\subsubsection{Scenarios}
\slidesrc{Terminology introduction: 14--15}

\begin{itemize}[nosep]
  \item \textbf{Closed MAS}: a known set of agents, e.g.\ the classic domestic-robot
        example with a \emph{robot}, an \emph{owner} and a \emph{supermarket} (the robot
        brings beer from the fridge to the owner and orders more from the supermarket).
        A user can ask ``Which are the available agents?'', ask the robot which plans it
        has, or ask it to describe a plan, and receives a natural-language answer.
  \item \textbf{Open MAS}: agents can join and leave, e.g.\ an auction of flight tickets
        where agents (Bob, Giacomo, Maria, \ldots) post bids such as ``I have a bid for the
        flight from Athens to Paris on 18/12/2015 for 50 euros'' and the winner is
        announced in natural language.
\end{itemize}

\subsubsection{Architecture (kept simple)}
\slidesrc{Terminology introduction: 16--17}

A human writes a sentence addressed to an agent, say ``@Bob, \ldots''. Before that, ChatBDI
has asked Bob for its \textbf{lexicon}: the list of terms (predicates) Bob understands.

\begin{center}
\begin{tikzpicture}[>=Stealth, node distance=0.55cm and 0.6cm,
    st/.style={draw, rounded corners, fill=nbgray!8, align=center, font=\footnotesize,
               minimum height=0.75cm},
    io/.style={draw, rounded corners, align=center, font=\footnotesize, minimum height=0.75cm}]
  \node[io, fill=nbgreen!15, draw=nbgreen] (in) {human:\\``@Bob sentence''};
  \node[st, right=1.1cm of in] (cls) {speech-act\\classifier};
  \node[st, above=of cls] (emb) {embedding of\\the sentence};
  \node[st, right=of cls] (sa) {speech act\\(performative)};
  \node[st, right=of emb] (fit) {best fit between\\sentence and\\Bob's lexicon};
  \node[io, fill=nbblue!12, draw=nbblue, right=1.0cm of fit] (lex) {Bob's\\lexicon};
  \node[st, below=0.9cm of sa, fill=nborange!15, draw=nborange] (llm) {LLM prompted with sentence,\\speech act and best-fitting terms};
  \node[io, fill=nbblue!12, draw=nbblue, right=0.7cm of llm] (out) {ACL message\\for Bob: \\performative + content};
  \draw[->] (in) -- (cls);
  \draw[->] (in.north) |- (emb);
  \draw[->] (cls) -- (sa);
  \draw[->] (emb) -- (fit);
  \draw[->] (lex) -- (fit);
  \draw[->] (sa) -- (llm);
  \draw[->] (fit.east) -- ++(0.3,0) coordinate (t) -- (t |- llm.north);
  \draw[->, thick] (llm) -- (out);
  \begin{scope}[on background layer]
    \node[draw=nbgray, dashed, rounded corners, inner sep=6pt,
          fit=(cls)(emb)(sa)(fit)(llm), label={[font=\scriptsize\bfseries, nbgray]below left:ChatBDI}] {};
  \end{scope}
\end{tikzpicture}
\captionof{figure}{ChatBDI translating a human sentence into an ACL message. The
performative comes from a classifier, the content from an LLM guided by the terms of
the recipient's lexicon that best fit the sentence.}
\end{center}

\begin{enumerate}[nosep]
  \item A \textbf{speech-act classifier} decides the performative: is the human
        informing (\texttt{tell}), asking (\texttt{askOne}, \texttt{askAll}),
        requesting an action (\texttt{achieve})?
  \item The sentence is turned into an \textbf{embedding} (a vector) and compared with
        Bob's lexicon to find the terms that \textbf{best fit} it.
  \item An \textbf{LLM is prompted} with the sentence, the speech act and the
        best-fitting terms, and produces the content of the message in ACL.
\end{enumerate}

\subsubsection{Experiments: it is a hard problem}
\slidesrc{Terminology introduction: 18--22}

Many LLMs were compared on the two sub-tasks: getting the \textbf{performative} right and
getting the \textbf{message content} right (the latter scored lower overall). To measure
how close a generated term is to the expected one, the experiments use the
\textbf{Levenshtein distance}: the minimum number of character insertions, deletions and
substitutions needed to turn one string into the other (the \emph{minimum edit
distance}, a topic of the course). Three examples show what can go wrong:

\begin{nbexample}{a partially correct translation}
\raggedright
Sentence: ``Employer 7 attend JS course till 1/11/2025''.
\begin{itemize}[nosep]
  \item Performative: \texttt{tell}, expected \texttt{tell}: correct.
  \item Term: \texttt{training\_attended(7, JS, 1/11/2025)}; expected
        \texttt{training\_attended(7, js, "1/11/2025")}.
  \item Levenshtein distance $=4$: \texttt{J}$\to$\texttt{j} and
        \texttt{S}$\to$\texttt{s} (2 substitutions), plus the two quotation marks
        around the date (2 insertions).
\end{itemize}
Almost right, but for a symbolic agent \texttt{JS} and \texttt{js} are different
constants and \texttt{1/11/2025} (an arithmetic expression!) is not the string
\texttt{"1/11/2025"}: the agent's plan will not match. Generating the content took about
15 seconds.
\end{nbexample}

\begin{nbexample}{right form, wrong meaning}
\raggedright
Sentence: ``beans are very good for you, they contain 24 grams of protein per 100
grams''.
\begin{itemize}[nosep]
  \item Performative: \texttt{tell}, correct.
  \item Term: \texttt{calories(beans, unit(24, gr), 100)}; expected
        \texttt{protein(beans, unit(100, gr), 24)}. Levenshtein distance $=13$.
\end{itemize}
The term is syntactically perfect, but it says something different: the predicate is
wrong (\emph{calories} instead of \emph{protein}) and the two quantities are swapped.
This is a \emph{semantic} error.
\end{nbexample}

\begin{nbexample}{the wrong speech act}
\raggedright
Sentence: ``Tell me the pressure of both front tires''.
\begin{itemize}[nosep]
  \item Performative: \texttt{tell}; expected \texttt{askAll}: \textbf{incorrect}.
  \item Term: \texttt{tire\_pressure(\_, 28)}; expected
        \texttt{tire\_pressure(front, \_)}. Levenshtein distance $=7$.
\end{itemize}
The sentence starts with the verb ``tell'', but the human is not \emph{informing} the
agent of anything: they are \emph{asking} for information, about all (``both'') front
tyres, hence \texttt{askAll}. The surface form of a sentence and the act it performs can
differ: this is a \emph{pragmatic} error.
\end{nbexample}

\begin{nbwhy}[the three errors]
The three examples are errors at the three levels of linguistic analysis of
Section~\ref{sec:three-levels}: wrong \emph{form} of a constant (syntax), wrong
\emph{meaning} of the fact (semantics), wrong \emph{act} performed by the sentence
(pragmatics). They also show why traditional NLP is still needed next to LLMs: to check
and repair the LLM's output (Section~\ref{sec:why-traditional-nlp}).
\end{nbwhy}

\begin{nbsummary}[what is NLP]
\begin{itemize}[nosep]
  \item NLP: computers processing human language; \textbf{NLU} (text $\to$
        representation) and \textbf{NLG} (representation $\to$ text).
  \item Eras: \textbf{symbolic} (1950s--early 1990s; ELIZA, SHRDLU, PARRY, MARGIE),
        \textbf{statistical} (1990s--2010s; corpora, probabilities, HMM POS tagging),
        \textbf{neural} (end-to-end, learnt features, word embeddings),
        \textbf{transformers} (BERT), \textbf{generative LLMs} (GPT, ChatGPT 2022).
  \item Traditional NLP is still needed when data are scarce, for pre/post-processing and
        validation of model output, and when LLMs cannot be used (technical, legal,
        ethical reasons).
  \item ChatBDI: BDI agents ``think'', an LLM ``talks''; a sentence becomes performative
        + content; errors appear at the syntactic, semantic and pragmatic level;
        Levenshtein distance measures how far the generated term is from the expected one.
\end{itemize}
\end{nbsummary}
